\section{שיעור 08 — חבורת $SO(N)$ ואלגברת $so(3)$ (16.11.2025)}

\subsection*{שאלות מנחות}
\begin{enumerate}
    \item מהי ההגדרה של חבורת הסיבובים $SO(N)$?
    \item כיצד מוצאים את הגנרטורים (היוצרים) של החבורה, המגדירים את "אלגברת לי" ($so(N)$) שלה?
    \item מהם יחסי החילוף (\LR{commutation relations}) של אלגברת $so(3)$ וכיצד הם קשורים לתנע זוויתי?
    \item כיצד ניתן, מתוך יחסי החילוף בלבד, למצוא את כל ההצגות הבלתי-פריקות (\LR{irreducible representations}) של תנע זוויתי?
    \item מה הקשר בין ההצגות הללו לבין $SO(3)$ ו-$SU(2)$?
\end{enumerate}

\subsection*{רעיונות מרכזיים והגדרות}

\begin{definitionbox}{חבורת הסיבובים $SO(N)$}
החבורה $SO(N)$ (\LR{Special Orthogonal group}) מוגדרת כאוסף כל המטריצות ה-$N \times N$ הממשיות, $R$, המקיימות שני תנאים:
\begin{enumerate}
    \item \textbf{אורתוגונליות:} $R^T R = I$ (כאשר $I$ היא מטריצת היחידה). תנאי זה מבטיח שהמטריצה משמרת אורך (נורמה).
    \item \textbf{"מיוחד" (\LR{Special}):} $\det(R) = 1$. תנאי זה מבטיח שהמטריצה מייצגת סיבוב "טהור" ללא שיקוף (שיהפוך את כיווניות המרחב).
\end{enumerate}
החבורה סגורה תחת כפל מטריצות, מכילה איבר יחידה, ולכל איבר יש הופכי ($R^{-1} = R^T$).
\end{definitionbox}

\begin{definitionbox}{אלגברת לי $so(N)$ (היוצרים)}
אלגברת לי $so(N)$ היא המרחב המשיק לחבורה $SO(N)$ בנקודת היחידה $I$.
אנו מוצאים אותה על ידי בחינת סיבובים קטנים (אינפיניטסימליים):
$$ R \approx I + \Omega $$
כאשר $\Omega$ היא מטריצה קטנה. נדרוש את תנאי האורתוגונליות:
\begin{align*}
I &= R^T R \\
I &\approx (I + \Omega)^T (I + \Omega) \\
I &\approx (I + \Omega^T) (I + \Omega) \\
I &\approx I + \Omega + \Omega^T + \mathcal{O}(\Omega^2)
\end{align*}
כדי ששוויון זה יתקיים, אנו חייבים לדרוש (בסדר ראשון):
$$ \Omega + \Omega^T = 0 \implies \Omega^T = -\Omega $$
\textbf{מסקנה:} אלגברת לי $so(N)$ היא אוסף כל המטריצות ה-$N \times N$ ה\textbf{אנטי-סימטריות} הממשיות.
\end{definitionbox}

\begin{notebox}{מספר הפרמטרים (דרגות החופש)}
מספר האיברים הבלתי תלויים במטריצה אנטי-סימטרית $N \times N$ הוא מספר האיברים שמעל (או מתחת) לאלכסון הראשי (שהוא אפס).
$$ \text{Degrees of Freedom} = \frac{N(N-1)}{2} $$
עבור $N=3$ (המרחב התלת-ממדי), אנו מקבלים $\frac{3(2)}{2} = 3$ פרמטרים, כצפוי (סיבוב סביב ציר $x, y, z$).
\end{notebox}

\subsection*{המעבר לאלגברת $so(3)$ ויחסי החילוף}
עבור $N=3$, יש לנו 3 יוצרים (\LR{generators}) בלתי תלויים, $J_1, J_2, J_3$, הפורסים את מרחב המטריצות האנטי-סימטריות $3 \times 3$.
(בהרצאה הוגדרו $X_{ab}$ באופן שקול).

לאחר חישוב יחסי החילוף בין מטריצות הבסיס (עם קונבנציה פיזיקלית הכוללת פקטור $i$ כדי להפוך את היוצרים להרמיטיים, כמקובל במכניקת קוונטים), מגיעים לתוצאה המרכזית:

\begin{formulabox}{אלגברת $so(3)$ (או $su(2)$)}
היוצרים של חבורת הסיבובים בתלת-ממד מקיימים את יחסי החילוף (האלגברה) של תנע זוויתי:
$$ [J_a, J_b] = i \epsilon_{abc} J_c $$
(כאשר $\hbar=1$ ו$\epsilon_{abc}$ הוא סמל לוי-צ'יויטה).
\end{formulabox}

\begin{importantbox}{משמעות התוצאה}
האופרטורים הדיפרנציאליים של תנע זוויתי ($L_x, L_y, L_z$) שמצאנו מפתרון מרחבי, והמטריצות $3 \times 3$ (הגנרטורים) של $SO(3)$, הן פשוט שתי \textbf{הצגות} שונות של אותה אלגברה אבסטרקטית המוגדרת על ידי יחסי החילוף לעיל.
כעת נראה שניתן למצוא את \textbf{כל} ההצגות האפשריות רק מתוך האלגברה.
\end{importantbox}

\subsection*{מציאת כל ההצגות הבלתי-פריקות (\LR{Irreps}) מן האלגברה}
אנו מתחילים אך ורק מההנחה שיש לנו שלושה אופרטורים (יוצרים) $J_1, J_2, J_3$ המקיימים $[J_a, J_b] = i \epsilon_{abc} J_c$.

\begin{enumerate}
    \item \textbf{הגדרת אופרטור קזימיר\LR{ (Casimir Operator):} }
    נגדיר את האופרטור $J^2 = J_1^2 + J_2^2 + J_3^2$.
    ניתן להראות (כפי שנעשה בהרצאה) שהוא חילופי עם כל היוצרים:
    $$ [J^2, J_a] = 0 \quad (\text{for } a=1,2,3) $$
    
    \item \textbf{בחירת בסיס משותף:}
    מכיוון ש-$J^2$ חילופי עם כל היוצרים, הוא בפרט חילופי עם $J_3$.
    $$ [J^2, J_3] = 0 $$
    לכן, ניתן למצוא בסיס של מצבים עצמיים משותפים לשני אופרטורים אלו. נסמן מצב בסיס ב-$|a, b\rangle$ (כאשר $a,b$ הם הערכים העצמיים):
    \begin{align*}
        J^2 |a, b\rangle &= a |a, b\rangle \\
        J_3 |a, b\rangle &= b |a, b\rangle
    \end{align*}

    \item \textbf{הגדרת אופרטורי סולם \LR{ (Ladder Operators):} }
    נגדיר בסיס חדש ("פולרי") לאלגברה: $J_0 \equiv J_3$ ו-
    $$ J_{\pm} = J_1 \pm i J_2 $$
    (נשים לב: $J_+^\dagger = J_-$ ו-$J_-^\dagger = J_+$ בהנחה ש-$J_1, J_2$ הרמיטיים).
    
    \item \textbf{יחסי חילוף בבסיס החדש:}
    מתוך יחסי החילוף המקוריים, ניתן לחשב:
    \begin{align*}
        [J_0, J_{\pm}] &= \pm J_{\pm} \\
        [J_+, J_-] &= 2J_0
    \end{align*}
    
    \item \textbf{פעולת אופרטורי הסולם על המצבים:}
    נבדוק כיצד $J_{\pm}$ פועלים על המצב $|a, b\rangle$:
    \begin{itemize}
        \item \textbf{פעולת $J^2$:} מכיוון ש-$[J^2, J_{\pm}] = 0$, נקבל:
        $$ J^2 (J_{\pm} |a, b\rangle) = J_{\pm} J^2 |a, b\rangle = a (J_{\pm} |a, b\rangle) $$
        \textbf{מסקנה:} $J_{\pm}$ אינם משנים את הערך העצמי $a$.
        
        \item \textbf{פעולת $J_0$:}
        \begin{align*}
            J_0 (J_{\pm} |a, b\rangle) &= ([J_0, J_{\pm}] + J_{\pm} J_0) |a, b\rangle \\
            &= (\pm J_{\pm} + J_{\pm} b) |a, b\rangle \\
            &= (b \pm 1) (J_{\pm} |a, b\rangle)
        \end{align*}
        \textbf{מסקנה:} $J_{\pm}$ מעלים או מורידים את הערך העצמי $b$ ב-1.
    \end{itemize}
    לכן: $J_{\pm} |a, b\rangle \propto |a, b \pm 1\rangle$.
    
    \item \textbf{הוכחת חסם \LR{ (Bounds on $b$):} }
    נשים לב כי $J_1^2 + J_2^2 = J^2 - J_0^2$.
    הערך העצמי של $J_1^2 + J_2^2$ חייב להיות אי-שלילי (כי אלו אופרטורים הרמיטיים בריבוע).
    $$ \langle a, b | (J^2 - J_0^2) | a, b \rangle = (a - b^2) \ge 0 $$
    \begin{importantbox}{תנאי חסם}
    עבור $a$ נתון, הערכים $b$ חייבים לקיים: $a \ge b^2$.
    \end{importantbox}
    
    \item \textbf{מצבי קצה \LR{(Maximum and Minimum States):} }
    הסולם חייב להיות סופי. לכן, חייב להתקיים:
    \begin{itemize}
        \item מצב עליון $b_{\text{max}}$: $J_+ |a, b_{\text{max}}\rangle = 0$
        \item מצב תחתון $b_{\text{min}}$: $J_- |a, b_{\text{min}}\rangle = 0$
    \end{itemize}
    
    \item \textbf{מציאת הערכים העצמיים המדויקים:}
    נשתמש בזהויות:
    $J_- J_+ = J^2 - J_0^2 - J_0$
    $J_+ J_- = J^2 - J_0^2 + J_0$
    
    נפעל על המצב העליון:
    $$ J_- J_+ |a, b_{\text{max}}\rangle = (J^2 - J_0^2 - J_0) |a, b_{\text{max}}\rangle = 0 $$
    $$ (a - b_{\text{max}}^2 - b_{\text{max}}) |a, b_{\text{max}}\rangle = 0 \implies \boxed{a = b_{\text{max}}(b_{\text{max}} + 1)} $$
    
    נפעל על המצב התחתון:
    $$ J_+ J_- |a, b_{\text{min}}\rangle = (J^2 - J_0^2 + J_0) |a, b_{\text{min}}\rangle = 0 $$
    $$ (a - b_{\text{min}}^2 + b_{\text{min}}) |a, b_{\text{min}}\rangle = 0 \implies \boxed{a = b_{\text{min}}(b_{\text{min}} - 1)} $$
    
    \item \textbf{סימטריה בין $b_{\text{max}}$ ו-$b_{\text{min}}$:}
    נשווה את שתי התוצאות עבור $a$:
    $$ b_{\text{max}}(b_{\text{max}} + 1) = b_{\text{min}}(b_{\text{min}} - 1) $$
    $$ (b_{\text{max}} + b_{\text{min}})(b_{\text{max}} - b_{\text{min}} + 1) = 0 $$
    מכיוון ש-$b_{\text{max}} \ge b_{\text{min}}$, הגורם השני חיובי, ולכן:
    $$ b_{\text{max}} + b_{\text{min}} = 0 \implies \boxed{b_{\text{min}} = -b_{\text{max}}} $$
    
    \item \textbf{כימות (\LR{Quantization}):}
    נסמן $j \equiv b_{\text{max}}$. לכן $b_{\text{min}} = -j$.
    הסולם מחבר את $-j$ ל-$j$ בצעדים של 1. מספר הצעדים $k$ חייב להיות שלם.
    $$ b_{\text{max}} - b_{\text{min}} = j - (-j) = 2j = k $$
    מכיוון ש-$k$ חייב להיות מספר שלם, $j$ יכול להיות:
    $$ j = 0, \frac{1}{2}, 1, \frac{3}{2}, 2, \dots $$
\end{enumerate}

\subsection*{סיכום: ההצגות הבלתי-פריקות של $so(3)$}
כל הצגה בלתי-פריקה (\LR{irrep}) של אלגברת $so(3)$ מאופיינת באופן ייחודי על ידי המספר $j$ (שלם או חצי-שלם).
נהוג לשנות את הסימון $a \to j(j+1)$ ו-$b \to m$.

\begin{formulabox}{ההצגות של $so(3) \cong su(2)$}
\begin{itemize}
    \item לכל הצגה יש מספר קוונטי $j \in \{0, \frac{1}{2}, 1, \frac{3}{2}, \dots \}$.
    \item מימד ההצגה (מספר המצבים) הוא $2j+1$.
    \item בסיס ההצגה מורכב ממצבים $|j, m\rangle$ כאשר $m \in \{-j, -j+1, \dots, j-1, j\}$.
    \item פעולת היוצרים על הבסיס:
    \begin{align*}
        J^2 |j, m\rangle &= j(j+1) |j, m\rangle \\
        J_z |j, m\rangle &= m |j, m\rangle \\
        J_{\pm} |j, m\rangle &= \sqrt{j(j+1) - m(m\pm 1)} |j, m \pm 1\rangle
    \end{align*}
\end{itemize}
\end{formulabox}

\subsection*{קשרים והמשך}

\begin{summarybox}{הקשר בין $SO(3)$ ל-$SU(2)$}
האלגברה שמצאנו, $[J_a, J_b] = i \epsilon_{abc} J_c$, היא אלגברת לי (\LR{Lie Algebra}) \textbf{זהה} עבור שתי חבורות לי (\LR{Lie Groups}) שונות: $SO(3)$ ו-$SU(2)$.
ההבדל ביניהן מתגלה כאשר מסתכלים על החבורה עצמה (סיבובים סופיים).

נבחן סיבוב ב-$2\pi$ סביב ציר $z$: $R_z(2\pi) = \exp(i 2\pi J_z)$.
$$ R_z(2\pi) |j, m\rangle = e^{i 2\pi m} |j, m\rangle $$
\begin{itemize}
    \item \textbf{עבור $SO(3)$:} סיבוב ב-$2\pi$ הוא בדיוק איבר היחידה ($I$). לכן $R_z(2\pi) = I$, שדורש $e^{i 2\pi m} = 1$. זה מתקיים רק אם $m$ הוא \textbf{מספר שלם}.
    \textbf{מסקנה:} ההצגות של $SO(3)$ הן אלו עם $j$ \textbf{שלם} ($j=l=0, 1, 2, \dots$). אלו ההצגות ה"בוזוניות", כמו הספירהרמוניות ($Y_{lm}$).
    
    \item \textbf{עבור $SU(2)$:} חבורה זו אינה דורשת שסיבוב ב-$2\pi$ יחזיר אותנו ליחידה (היא "כיסוי כפול" של $SO(3)$).
    כאשר $j$ \textbf{חצי-שלם} ($j=1/2, 3/2, \dots$), $m$ הוא גם חצי-שלם, ונקבל $e^{i 2\pi m} = e^{i \pi (2m)} = -1$.
    \textbf{מסקנה:} ההצגות החצי-שלמות ("פרמיוניות", כמו ספין) הן הצגות "אמיתיות" של $SU(2)$, אך לא של $SO(3)$.
\end{itemize}
האלגברה $so(3) \cong su(2)$ (האלגברות איזומורפיות) מכילה את \textbf{כל} ההצגות האפשריות.
\end{summarybox}

\subsection*{הבהרות נדרשות \LR{ (Clarifications Needed)}}
\begin{warningbox}{בלבול בין $SO(N)$ ל-$SU(N)$ בהרצאה}
במהלך הגזירה של יחסי החילוף עבור $SO(N)$, ההרצאה השתמשה בקונבנציה (הגדרת $X_{ab}$ עם פקטור $i$) שהופכת את היוצרים להרמיטיים. יוצרים הרמיטיים (כאשר מוכפלים ב-$i$ באקספוננט) יוצרים חבורה \textbf{אוניטרית} ($SU(N)$), לא אורתוגונלית ($SO(N)$).
היוצרים של $SO(N)$ הם מטריצות ממשיות אנטי-סימטריות (שהן אנטי-הרמיטיות).
עם זאת, התוצאה הסופית עבור $N=3$ היא שיחסי החילוף זהים לאלו של $su(2)$, כלומר $so(3) \cong su(2)$, ולכן הניתוח האלגברי של ההצגות נותר נכון ותקף.
\end{warningbox}


\section*{הרצאה 8: סימטריות, תנע זוויתי וחבורת $SO(3)$ במכניקה קוונטית}

\subsection*{מטרות ההרצאה}
בהרצאה זו נעמיק בהבנת הקשר בין סימטריות פיזיקליות לבין תנע זוויתי קוונטי. נבחן את מבנה האלגברה של $so(3)$ ושל הכיסוי הכפול שלה $su(2)$, נלמד כיצד נבנות ההצגות הבלתי-פריקות שלה, ונראה כיצד הן מתבטאות במצבי תנע זוויתי כמו הפונקציות הספירהרמוניות $Y_{lm}$.

---

\begin{importantbox}
\textbf{רעיון מרכזי:}
במכניקה קוונטית, לכל סימטריה יש גודל שמור. במקרה של סימטריה סיבובית, הגודל השמור הוא התנע הזוויתי.
\end{importantbox}

---

\subsection*{1. הגדרה: תנע זוויתי קוונטי}

נגדיר את אופרטורי התנע הזוויתי:

\[
\vec{J} = (J_x, J_y, J_z)
\]

הם מקיימים את יחסי הקומוטציה:
\[
[J_i, J_j] = i\hbar \, \epsilon_{ijk} J_k
\]

ובפרט:
\[
[J_z, J_{\pm}] = \pm \hbar J_{\pm}, \quad [J_+, J_-] = 2\hbar J_z
\]

כאשר \( J_{\pm} = J_x \pm i J_y \).

---

\subsection*{2. הערכים העצמיים של $\vec{J}^2$ ושל $J_z$}

נניח כי קיימים וקטורים עצמיים משותפים לשני האופרטורים הקומוטטיביים $\{J^2, J_z\}$:
\[
J^2 |j, m\rangle = \hbar^2 j(j+1) |j, m\rangle, \quad J_z |j, m\rangle = \hbar m |j, m\rangle
\]

כאשר:
\[
m = -j, -j+1, \ldots, j-1, j
\]
והמימד הכולל של ההצגה הוא:
\[
\text{dim} = 2j + 1
\]

\begin{examplebox}{פעולת אופרטורי העלאה והורדה}
ניישם את $J_\pm$ על בסיס המצבים:
\[
J_{\pm}|j,m\rangle = \hbar \sqrt{j(j+1) - m(m\pm1)}\, |j, m\pm1\rangle
\]

הפעולה מעלה או מורידה את ערך $m$ ב-1, אך אינה משנה את $j$.
\end{examplebox}

---

\subsection*{3. מצבים עליון ותחתון}

נגדיר:
\[
J_+ |j, m_{\text{max}}\rangle = 0, \quad J_- |j, m_{\text{min}}\rangle = 0
\]
מתוך יחסי הקומוטציה נובע:
\[
a = m_{\text{max}}(m_{\text{max}} + 1) = m_{\text{min}}(m_{\text{min}} - 1)
\]
ומכאן:
\[
m_{\text{min}} = -m_{\text{max}} \equiv -j
\]
כלומר, הערכים של $m$ פרוסים מ־$-j$ ועד $+j$ בקפיצות של 1.

---

\subsection*{4. חבורות $SO(3)$ ו־$SU(2)$}

נבחן סיבוב מלא ב-$2\pi$ סביב ציר $z$:
\[
R_z(2\pi) = e^{i 2\pi J_z / \hbar}
\]
כאשר פועלים על מצב $|j,m\rangle$ נקבל:
\[
R_z(2\pi)|j,m\rangle = e^{i 2\pi m}|j,m\rangle
\]

\begin{itemize}
    \item עבור \textbf{$SO(3)$}: סיבוב ב-$2\pi$ הוא זהות $\Rightarrow e^{i2\pi m}=1 \Rightarrow m$ שלם.
    \item עבור \textbf{$SU(2)$}: סיבוב ב-$2\pi$ אינו זהות; שהוא "כיסוי כפול" של $SO(3)$. לכן $m$ יכול להיות גם חצי־שלם.
\end{itemize}

\begin{formulabox}{מסקנה}
ל־$SO(3)$ יש הצגות רק עבור $j$ שלמים (בוזונים),  
ואילו ל־$SU(2)$ — גם חצי־שלמים (פרמיונים).
\end{formulabox}

---

\subsection*{5. דוגמה פיזיקלית: הפונקציות הספירהרמוניות}

הפתרונות של משוואת שרדינגר לחלקיק בקואורדינטות כדוריות כוללים גורם זוויתי מהצורה \( Y_{lm}(\theta,\phi) \), המקיים:
\[
L^2 Y_{lm} = \hbar^2 l(l+1) Y_{lm}, \quad L_z Y_{lm} = \hbar m Y_{lm}
\]
כאשר $l$ שלם בלבד.  
מכאן שה־$Y_{lm}$ מהווים בדיוק את ההצגות של $SO(3)$.

---

\subsection*{6. סימטריה נסתרת במערכת המימן}

עבור האטום ההידרוגני, האנרגיה תלויה רק במספר הקוונטי הראשי $n$ ולא ב-$l$.  
הניוון ברמה הוא $n^2$, מה שמרמז על סימטריה גדולה יותר מ־$SO(3)$.

\begin{theorembox}{וקטור רונגה–לנצ}
מוגדר:
\[
\vec{A} = \frac{1}{2\mu}(\vec{p}\times\vec{L} - \vec{L}\times\vec{p}) - \frac{Ze^2}{4\pi\epsilon_0}\frac{\vec{r}}{r}
\]
וקטור זה מקיים $[\hat{H}, \vec{A}] = 0$, ולכן הוא גודל שמור.
הגדלים $\vec{L}$ ו־$\vec{A}$ יוצרים יחד את החבורה $SO(4)$, המסבירה את הניוון $n^2$.
\end{theorembox}

\begin{table}[h!]
\centering
\begin{tabular}{|c|c|c|c|c|}
\hline
\textbf{חבורה} & \textbf{אלגברה} & \textbf{ערכים מותרים של $j$} & \textbf{ממד ההצגה} & \textbf{סיבוב ב-$2\pi$} \\
\hline
$SO(3)$ & $so(3)$ & $j = 0, 1, 2, \dots$ & $2j + 1$ & זהות ($R_z(2\pi)=I$) \\
\hline
$SU(2)$ & $su(2)$ & $j = 0, \tfrac{1}{2}, 1, \tfrac{3}{2}, \dots$ & $2j + 1$ & $R_z(2\pi) = (-1)^{2j}$ \\
\hline
\end{tabular}
\caption{השוואה בין ההצגות של $SO(3)$ ושל $SU(2)$ \cite[id=1]{}.}
\end{table}


\begin{center}
\begin{tikzpicture}[scale=1.2, thick]
% Circles
\draw[blue!60, very thick] (0,0) circle (1.5);
\draw[green!60, very thick, dashed] (0,2) circle (1.5);

% Arrows showing covering
\draw[->, thick] (0,0.5) -- (0,1.5);
\node[right] at (0,1) {כיסוי כפול};

% Labels
\node at (0,-2.1) {חבורה $SO(3)$};
\node at (0,3.7) {חבורה $SU(2)$};
\end{tikzpicture}
\end{center}

\subsection*{7. סיכום ורעיונות מרכזיים}
\begin{itemize}
    \item חבורת הסיבובים $SO(3)$ מתוארת ע"י האלגברה $so(3)$, איזומורפית ל־$su(2)$.
    \item ההצגות מתויגות ע"י מספר קוונטי $j$ ושלושה אופרטורים $J_x, J_y, J_z$.
    \item עבור $SO(3)$ — רק $j$ שלם; עבור $SU(2)$ — גם $j$ חצי־שלם.
    \item הפונקציות $Y_{lm}$ הן בסיס להצגות של $SO(3)$.
    \item באטום המימן מופיעה סימטריה נסתרת $SO(4)$ שגורמת לניוון הגבוה $n^2$ [3].
\end{itemize}

\begin{summarybox}
\textbf{תובנה מרכזית:}  
הבנת חבורות הסימטריה במכניקה קוונטית אינה רק עניין מתמטי –  
היא מגלה מבנים שמורים ופיזיקליים עמוקים כמו הניוון באטום המימן.
\end{summarybox}